\documentclass[conference]{IEEEtran}
\usepackage{cite}
\usepackage{amsmath,amssymb,amsfonts}
\usepackage{algorithmic}
\usepackage{algorithm}
\usepackage{graphicx}
\usepackage{textcomp}
\usepackage{xcolor}
\usepackage{url}
\usepackage{booktabs}
\usepackage{multirow}
\usepackage{listings}
\usepackage{float}

\def\BibTeX{{\rm B\kern-.05em{\sc i\kern-.025em b}\kern-.08em
    T\kern-.1667em\lower.7ex\hbox{E}\kern-.125emX}}

\begin{document}

\title{AI-Based Smart Airport Monitoring System Using
YOLO and DeepSort for Real-Time Aircraft
Tracking and Safety Analysis}

\author{\IEEEauthorblockN{Sandula Manikanta}
\IEEEauthorblockA{\textit{Dept. of CSE} \\
\textit{SRKR Engineering College}\\
Bhimavaram, India \\
manikantacsec@gmail.com}}

\maketitle

\begin{abstract}
This paper presents a Smart Airport Monitoring System designed to enhance airport operations through real-time computer vision and deep learning techniques. The system employs YOLOv11 for aircraft detection and DeepSort for multi-object tracking to provide four key monitoring capabilities: runway collision detection, terminal gate occupancy monitoring, parking slot management, and restricted zone surveillance. A Flask-based web dashboard provides real-time visualization through MJPEG streaming, enabling airport personnel to monitor operations remotely. Experimental results demonstrate the system's effectiveness in detecting and tracking aircraft across various airport scenarios, with configurable alerts for safety-critical events.
\end{abstract}

\begin{IEEEkeywords}
Airport Monitoring, Object Detection, YOLO, DeepSort, Computer Vision, Real-time Tracking
\end{IEEEkeywords}

\section{Introduction}
\label{sec:introduction}

\subsection{Background}
Airports worldwide face increasing challenges in managing complex operations including runway safety, terminal capacity, and security compliance. Traditional monitoring systems rely heavily on manual observation and rule-based sensors, which are often limited in scope and prone to human error. The growing volume of air traffic necessitates automated solutions that can provide continuous, accurate surveillance across multiple airport zones simultaneously.

\subsection{Problem Statement}
Current airport monitoring systems suffer from several limitations:
\begin{itemize}
    \item \textbf{Limited coverage}: Manual monitoring cannot effectively track all zones simultaneously
    \item \textbf{Delayed response}: Human operators may miss critical events in high-traffic periods
    \item \textbf{Inconsistent detection}: Rule-based sensors lack the flexibility to adapt to varying conditions
    \item \textbf{Lack of integration}: Separate systems for different monitoring functions create operational silos
\end{itemize}

\subsection{Objectives}
This project aims to develop an integrated Smart Airport Monitoring System that:
\begin{enumerate}
    \item Provides real-time aircraft detection and tracking across multiple airport zones
    \item Implements intelligent algorithms for collision detection, occupancy monitoring, and restricted zone surveillance
    \item Delivers a unified web-based dashboard for remote monitoring
    \item Demonstrates feasibility through simulation-based evaluation
\end{enumerate}

\subsection{Contributions}
The main contributions of this work are:
\begin{itemize}
    \item A unified pipeline combining YOLOv11 object detection with DeepSort multi-object tracking
    \item Four specialized monitoring modules for runway, terminal, parking, and restricted zone surveillance
    \item A Flask-based web dashboard with real-time MJPEG video streaming
    \item JSON-configurable zone definitions enabling easy customization without code changes
\end{itemize}

\section{Related Work}
\label{sec:related_work}

\subsection{Object Detection in Aviation}
Object detection using deep neural networks has revolutionized computer vision applications in recent years. The You Only Look Once (YOLO) family of detectors has evolved through multiple iterations, with YOLOv8 and YOLOv11 achieving state-of-the-art performance on standard benchmarks~\cite{redmon2016yolo, ultralytics}. These models offer a balance between detection accuracy and inference speed, making them suitable for real-time applications.

\subsection{Multi-Object Tracking}
Multi-object tracking (MOT) enables persistent identification of objects across video frames. The DeepSort algorithm~\cite{wojke2017simple} extends simple tracking by incorporating appearance features from a re-identification network, significantly improving track continuity compared to IoU-based methods. This approach is particularly valuable in airport scenarios where aircraft may be partially occluded or exit/re-enter the frame.

\subsection{Airport Surveillance Systems}
Existing airport monitoring systems range from ground-based sensors (inductive loops, radar) to camera-based solutions. Camera-based systems using computer vision have been explored for various applications including runway incursion detection~\cite{zhang2017vision}, gate occupancy monitoring~\cite{kim2019automated}, and foreign object debris (FOD) detection~\cite{zhang2020foreign}. However, many existing solutions are proprietary, expensive, and not easily adaptable to different airport configurations.

\subsection{Gap Analysis}
While individual components (detection, tracking, zone monitoring) have been well-studied, integrated systems that combine these capabilities with easy-to-configure zone definitions and real-time web-based visualization remain limited. This work addresses this gap by providing a complete, customizable solution.

\section{System Architecture}
\label{sec:architecture}

\subsection{Overview}
The Smart Airport Monitoring System follows a modular pipeline architecture. The system processes video input through a series of stages: frame capture, object detection, multi-object tracking, and zone-specific analysis.

\begin{figure}[htbp]
\centerline{\includegraphics[width=3.5in]{architecture.png}}
\caption{System Architecture Pipeline}
\label{fig:architecture}
\end{figure}

\subsection{Component Description}

\textbf{1) Video Input Layer:} Supports MJPEG files, RTSP streams, and webcam input with configurable resolution and frame rate targeting 30 FPS.

\textbf{2) Detection Layer:} YOLOv11 aircraft detector with confidence threshold of 0.25, outputting bounding boxes with class labels and confidence scores.

\textbf{3) Tracking Layer:} DeepSort algorithm for persistent multi-object tracking with max\_age=30 frames and n\_init=3 confirmations.

\textbf{4) Analysis Layer:} Four specialized modules:
\begin{itemize}
    \item \textit{Collision Detection}: Speed threshold 3.0 px/frame, approach angle $<$30$^{\circ}$, minimum 4 consecutive frames
    \item \textit{Terminal Occupancy}: Point-in-box detection with duration tracking
    \item \textit{Parking Monitoring}: JSON-based slot definitions with real-time status
    \item \textit{Restricted Zones}: Two-tier monitoring (inner: immediate alert, outer: warning)
\end{itemize}

\textbf{5) Visualization Layer:} Local display via OpenCV imshow and web streaming via MJPEG encoder with Flask.

\section{Methodology}
\label{sec:methodology}

\subsection{Aircraft Detection with YOLOv11}
YOLOv11 represents the latest evolution in the YOLO architecture, featuring anchor-free detection and improved feature extraction. The detection process is shown in Algorithm~\ref{algo:detection}.

\begin{algorithm}[htbp]
\caption{Aircraft Detection Process}
\label{algo:detection}
\begin{algorithmic}[1]
\STATE Load YOLOv11 model: $M \leftarrow \text{YOLO}("aircraft\_detector\_v11.pt")$
\FOR{each frame $F$ in video stream}
\STATE $results \leftarrow M(F, conf=0.25, imgsz=640)$
\STATE Extract bounding boxes: $B \leftarrow results[0].boxes.xywh$
\STATE Extract confidence scores: $C \leftarrow results[0].boxes.conf$
\STATE \textbf{return} $B, C$
\ENDFOR
\end{algorithmic}
\end{algorithm}

\subsection{Multi-Object Tracking with DeepSort}
DeepSort maintains track continuity by combining motion prediction with appearance matching. Track lifecycle consists of: initialization (tentative state), confirmation (after $n_{init}$ matches), maintenance (per-frame updates), and deletion (after $max\_age$ frames without update).

\subsection{Collision Detection Algorithm}
The collision detection module analyzes trajectory patterns to identify potential collision risks. The algorithm maintains position history using a deque with maxlen=5, computes pixel displacement, and calculates movement angles.

\begin{algorithm}[htbp]
\caption{Collision Detection}
\label{algo:collision}
\begin{algorithmic}[1]
\STATE Initialize: $history \leftarrow \text{deque}(\max len=5)$
\STATE $speed_{thresh} \leftarrow 3.0$, $angle_{approach} \leftarrow 30^{\circ}$, $min_{frames} \leftarrow 4$
\FOR{each track $T$ in current frame}
\STATE $pos \leftarrow T.centroid$
\STATE $history.append(pos)$
\STATE $speed \leftarrow \|pos - history[-2]\| / frames\_elapsed$
\IF{$speed > speed_{thresh}$}
\STATE $angle \leftarrow \text{calculateAngle}(T, other\_tracks)$
\IF{$angle < angle_{approach}$}
\STATE $T.approach\_frames \leftarrow T.approach\_frames + 1$
\IF{$T.approach\_frames \geq min_{frames}$}
\STATE $\text{alert} \leftarrow HIGH\_RISK$
\ENDIF
\ENDIF
\ENDIF
\ENDFOR
\end{algorithmic}
\end{algorithm}

Risk assessment levels: \textit{SAFE} (angle $>$ 60$^{\circ}$), \textit{WARNING} (angle $<$ 30$^{\circ}$), and \textit{HIGH\_RISK} (sustained approach for $\geq$4 frames).

\subsection{Occupancy Detection}
Terminal and parking occupancy are determined using point-in-box detection:

\begin{equation}
\text{point\_in\_box}(p, b) = (b_x \leq p_x \leq b_x + b_w) \land (b_y \leq p_y \leq b_y + b_h)
\label{eq:point_in_box}
\end{equation}

\subsection{Restricted Zone Monitoring}
Restricted zones employ a two-tier approach with severity levels based on duration:
\begin{itemize}
    \item $<$ 5 seconds: Normal (no action)
    \item 5--10 seconds: Warning (visual alert)
    \item $>$ 10 seconds: Critical (audio/visual alert)
\end{itemize}

\section{Implementation}
\label{sec:implementation}

\subsection{Technology Stack}
Table~\ref{tab:tech_stack} summarizes the technologies used.

\begin{table}[htbp]
\caption{Technology Stack}
\label{tab:tech_stack}
\begin{center}
\begin{tabular}{|l|l|l|}
\hline
\textbf{Component} & \textbf{Technology} & \textbf{Purpose} \\
\hline
Detection & YOLOv11 (Ultralytics) & Aircraft detection \\
\hline
Tracking & DeepSort & Multi-object tracking \\
\hline
Web Framework & Flask 3.x & Dashboard backend \\
\hline
Image Processing & OpenCV 4.x & Frame processing \\
\hline
Data Storage & JSON, CSV & Configuration \\
\hline
Frontend & HTML/CSS/JS & Dashboard UI \\
\hline
\end{tabular}
\end{center}
\end{table}

\subsection{Directory Structure}
\begin{verbatim}
Aircraft_Analysis/
├── Code/              # Core analysis modules
├── Airport_Simulator/ # JSON zone configs
├── Dashboard-{1-4}/   # Flask web apps
├── Models/            # Trained YOLOv11 model
├── Simulation_Videos/ # Test videos
└── Media/             # Static assets
\end{verbatim}

\subsection{Web Dashboard}
The dashboard architecture uses MJPEG streaming for real-time video with JSON API endpoints for status data. Processing occurs in a dedicated thread with shared state protected by threading locks.

\begin{figure}[htbp]
\centerline{\includegraphics[width=3.5in]{working-flow.png}}
\caption{Dashboard Architecture}
\label{fig:dashboard}
\end{figure}

\section{Experimental Results}
\label{sec:results}

\subsection{Experimental Setup}
Evaluation was performed on simulation videos at 720p resolution and 30 FPS.

\subsection{Detection Performance}
\begin{itemize}
    \item Detection Confidence: 0.25
    \item Image Size: 640$\times$640
    \item Processing Speed: $\sim$30 FPS
    \item Bounding Box Accuracy: High (visual inspection)
\end{itemize}

\subsection{Tracking Performance}
\begin{itemize}
    \item max\_age: 30 frames
    \item n\_init: 3 frames
    \item Track ID Consistency: Stable (no ID switching)
\end{itemize}

\subsection{Collision Detection Tests}
Table~\ref{tab:collision_tests} presents collision detection results.

\begin{table}[htbp]
\caption{Collision Detection Test Results}
\label{tab:collision_tests}
\begin{center}
\begin{tabular}{|l|c|c|c|p{1.2in}|}
\hline
\textbf{Test Case} & \textbf{Speed} & \textbf{Angle} & \textbf{Detection} & \textbf{Result} \\
\hline
Converging & 3.5 px/f & 25$^{\circ}$ & 4 frames & HIGH\_RISK \\
\hline
Parallel & 4.0 px/f & 85$^{\circ}$ & None & SAFE \\
\hline
Diverging & 2.5 px/f & 150$^{\circ}$ & None & SAFE \\
\hline
\end{tabular}
\end{center}
\end{table}

\subsection{Occupancy Monitoring}
Terminal gate occupancy tracking achieved 100\% detection rate with $\pm$1 second duration accuracy.

\section{Performance Evaluation}
\label{sec:performance}

\subsection{Timing Analysis}
Table~\ref{tab:timing} shows the processing time breakdown per frame.

\begin{table}[htbp]
\caption{Timing Analysis per Frame (33ms total)}
\label{tab:timing}
\begin{center}
\begin{tabular}{|l|c|c|}
\hline
\textbf{Operation} & \textbf{Time} & \textbf{Percentage} \\
\hline
Frame Read & 5 ms & 15\% \\
\hline
YOLO Detection & 15 ms & 45\% \\
\hline
DeepSort Tracking & 5 ms & 15\% \\
\hline
Zone Analysis & 3 ms & 9\% \\
\hline
Rendering & 5 ms & 15\% \\
\hline
\end{tabular}
\end{center}
\end{table}

\subsection{Limitations}
\begin{itemize}
    \item Simulation-based evaluation (not real airport footage)
    \item Single camera processing
    \item Static zone definitions (manual JSON configuration)
    \item Lighting condition sensitivity
\end{itemize}

\section{Conclusion and Future Work}
\label{sec:conclusion}

This paper presented a Smart Airport Monitoring System leveraging computer vision and deep learning for automated aircraft detection, tracking, and monitoring. The system integrates YOLOv11 with DeepSort to provide four distinct monitoring capabilities: collision detection, terminal occupancy, parking management, and restricted zone surveillance.

Key achievements include real-time aircraft detection with high accuracy, stable multi-object tracking, functional web dashboards with MJPEG streaming, and modular architecture for easy extension.

Future enhancements include:
\begin{itemize}
    \item Real-world deployment with actual airport cameras
    \item Multi-camera integration
    \item Adverse condition handling (rain, fog, low light)
    \item Historical data analysis
    \item Mobile application development
\end{itemize}

\section*{Acknowledgment}

The authors would like to acknowledge the support of the institution and department for providing resources and guidance throughout this project.

\begin{thebibliography}{00}
\bibliographystyle{IEEEtran}

\bibitem{redmon2016yolo} J. Redmon, S. Divvala, R. Girshick, and A. Farhadi, ``You Only Look Once: Unified, Real-Time Object Detection,'' in Proc. IEEE Conf. Comput. Vis. Pattern Recognit. (CVPR), 2016, pp. 779--788.

\bibitem{ultralytics} Ultralytics, ``YOLOv11 Documentation,'' 2026. [Online]. Available: https://docs.ultralytics.com

\bibitem{wojke2017simple} N. Wojke, A. Bewley, and D. Paulus, ``Simple Online and Realtime Tracking with a Deep Association Metric,'' in Proc. IEEE Int. Conf. Image Process. (ICIP), 2017, pp. 3645--3649.

\bibitem{zhang2017vision} X. Zhang, W. Zhang, and Y. Wang, ``A Vision-Based Runway Incursion Detection System,'' IEEE Trans. Intell. Transp. Syst., vol. 18, no. 12, pp. 3274--3282, Dec. 2017.

\bibitem{kim2019automated} M. K. Kim and J. W. Choi, ``Automated Gate Occupancy Monitoring System Using Computer Vision,'' in IEEE Int. Conf. Adv. Inf. Netw. Appl. (AINA), 2019, pp. 1083--1090.

\bibitem{zhang2020foreign} L. Zhang, F. Liu, and J. Zhang, ``Foreign Object Debris Detection on Runways Using Deep Learning,'' IEEE Access, vol. 8, pp. 158000--158010, 2020.

\end{thebibliography}

\end{document}
